\documentclass[12pt,a4paper]{article}
\usepackage[utf8]{inputenc}
\usepackage[french]{babel}
\usepackage[T1]{fontenc}
\usepackage{graphicx}
\usepackage{tikz}
\usepackage{listings}
\usepackage{xcolor}
\usepackage{hyperref}
\usepackage{geometry}
\usepackage{float}
\usepackage{titlesec}

% TikZ Libraries for UML
\usetikzlibrary{shapes, arrows.meta, positioning, shadows, fit, calc, backgrounds}

\geometry{margin=2.5cm}

% Code Style
\definecolor{codegreen}{rgb}{0,0.6,0}
\definecolor{codegray}{rgb}{0.5,0.5,0.5}
\definecolor{codepurple}{rgb}{0.58,0,0.82}
\definecolor{backcolour}{rgb}{0.95,0.95,0.92}

\lstdefinestyle{estiloJava}{
    backgroundcolor=\color{backcolour},   
    commentstyle=\color{codegreen},
    keywordstyle=\color{magenta},
    numberstyle=\tiny\color{codegray},
    stringstyle=\color{codepurple},
    basicstyle=\ttfamily\footnotesize,
    breakatwhitespace=false,         
    breaklines=true,                 
    captionpos=b,                    
    keepspaces=true,                 
    numbers=left,                    
    numbersep=5pt,                  
    showspaces=false,                
    showstringspaces=false,
    showtabs=false,                  
    tabsize=2,
    frame=single
}
\lstset{style=estiloJava}

% Title Format
\titleformat{\section}
  {\normalfont\Large\bfseries}{\thesection}{1em}{}

% Custom UML Definitions for TikZ (Simulating tikz-uml for portability)
\tikzset{
    actor/.style={
        circle, fill=white, draw=black, minimum size=0.5cm, inner sep=0pt,
        append after command={
            \pgfextra{
                \draw (\tikzlastnode.south) -- ++(0,-0.5) coordinate (body);
                \draw (body) -- ++(-0.3,-0.4);
                \draw (body) -- ++(0.3,-0.4);
                \draw (body) ++(0,0.3) -- ++(-0.4,0);
                \draw (body) ++(0,0.3) -- ++(0.4,0);
            }
        }
    },
    usecase/.style={
        ellipse, draw=black, fill=blue!10, text width=2.5cm, align=center, minimum height=1cm
    },
    class/.style={
        rectangle, draw=black, fill=yellow!10, text width=4cm,
        minimum height=0.5cm, node distance=0.5cm, font=\bfseries
    },
    attribute/.style={
        rectangle, draw=black, fill=white, text width=4cm,
        minimum height=0.5cm, node distance=0cm, align=left, font=\footnotesize
    },
    package/.style={
        rectangle, draw=black, fill=gray!5, minimum width=3cm, minimum height=2cm,
        inner sep=0.5cm, inner ysep=0.8cm
    },
    component/.style={
        rectangle, draw=black, fill=white, minimum width=3cm, minimum height=1.5cm,
        append after command={
            \pgfextra{
                \draw (\tikzlastnode.west) ++(-0.2, 0.4) rectangle ++(0.4, 0.2);
                \draw (\tikzlastnode.west) ++(-0.2, -0.2) rectangle ++(0.4, 0.2);
            }
        }
    }
}

\begin{document}

%--- TITLE PAGE ---
\begin{titlepage}
    \begin{center}
        \Large \textbf{Université Blida 1}\\
        \large Faculté des Sciences\\
        Département d'Informatique\\
        
        \vspace{2cm}
        \includegraphics[width=4cm]{example-image-a} % Placeholder for Logo
        
        \vspace{1.5cm}
        {\Huge \textbf{Rapport de Projet Architecture Logicielle}}\\[0.5cm]
        {\Large \textbf{Plateforme Oudekenis}}\\[0.5cm]
        
        \vspace{2cm}
        \textbf{Réalisé par:}\\
        Etudiant 1\\
        Etudiant 2\\
        
        \vspace{1cm}
        \textbf{Module:}\\
        Architecture Logicielle (Master 2 IL)\\
        
        \vspace{2cm}
        \textbf{Année Universitaire:}\\
        2025/2026
        
    \end{center}
    \vfill
    \centering
    Janvier 2026
\end{titlepage}

\tableofcontents
\newpage

%--- 1. INTRODUCTION ---
\section{Inroduction}
Ce projet s'inscrit dans le cadre du module d'Architecture Logicielle et vise à concevoir et développer une plateforme de commerce électronique similaire à \textit{Ouedkniss}, leader des annonces en Algérie. L'objectif est de mettre en pratique les concepts architecturaux modernes tels que le \textbf{Modular Monolith} (via Spring Modulith), l'intégration de moteurs de recherche avancés, et la communication temps réel.

\subsection{Analyse de l'existant (Ouedkniss)}
Le site Ouedkniss est une référence, mais il présente certaines limitations technologiques perçues :
\begin{itemize}
    \item \textbf{Navigation} : Parfois complexe avec une hiérarchie de catégories profonde.
    \item \textbf{Temps Réel} : La messagerie manque souvent d'instantanéité fluide.
    \item \textbf{Recherche} : La pertinence des résultats peut être améliorée (flou/typos).
\end{itemize}

\subsection{Notre Solution : Oudekenis}
Notre solution \textit{Oudekenis} propose :
\begin{itemize}
    \item Une architecture **Spring Modulith** pour une maintenabilité accrue.
    \item Un moteur de recherche **Elasticsearch** pour une indexation rapide et tolérante aux fautes.
    \item Une messagerie instantanée basée sur **WebSocket (STOMP)** pour relier acheteurs et vendeurs en temps réel.
    \item Une interface réactive avec **Thymeleaf** et **Alpine.js**.
\end{itemize}

\newpage

%--- 2. ANALYSE DES BESOINS (USE CASE) ---
\section{Analyse des Cas d'Utilisation}

Le diagramme ci-dessous illustre les interactions des principaux acteurs avec le système.

\begin{center}
\begin{tikzpicture}[scale=0.8, every node/.style={transform shape}]
    % Actors
    \node[actor] (Guest) at (0, 8) {}; \node[below] at (Guest) {Visiteur};
    \node[actor] (User) at (0, 4) {}; \node[below] at (User) {Utilisateur};
    \node[actor] (Seller) at (0, 0) {}; \node[below] at (Seller) {Vendeur};

    % System Box
    \draw (3, -2) rectangle (12, 10);
    \node[anchor=north] at (7.5, 10) {\textbf{Système Oudekenis}};

    % Use Cases
    \node[usecase] (Search) at (5, 9) {Rechercher Produit};
    \node[usecase] (View) at (5, 7.5) {Consulter Détail};
    \node[usecase] (Auth) at (7, 6) {S'authentifier};
    
    \node[usecase] (Chat) at (9, 4.5) {Chat Temps-Réel};
    \node[usecase] (Profile) at (9, 3) {Gérer Profil};
    \node[usecase] (Fav) at (5, 4.5) {Ajouter Favoris};
    
    \node[usecase] (Post) at (6, 1) {Publier Annonce};
    \node[usecase] (Update) at (6, -1) {Modifier Annonce};

    % Relationships
    \draw (Guest) -- (Search);
    \draw (Guest) -- (View);
    \draw (Guest) -- (Auth);

    % Inheritance
    \draw[-{Triangle[open]}] (User) -- (Guest);
    \draw[-{Triangle[open]}] (Seller) -- (User);

    % User Relations
    \draw (User) -- (Chat);
    \draw (User) -- (Profile);
    \draw (User) -- (Fav);

    % Seller Relations
    \draw (Seller) -- (Post);
    \draw (Seller) -- (Update);
    \draw[dashed, ->] (Post) -- node[right, font=\tiny] {<<include>>} (View);

\end{tikzpicture}
\end{center}

\subsection{Acteurs et Rôles}
\begin{itemize}
    \item \textbf{Visiteur} : Peut rechercher des produits via Elasticsearch et consulter les détails.
    \item \textbf{Utilisateur (Connecté)} : Peut gérer son profil, ajouter des favoris et discuter avec les vendeurs via WebSocket.
    \item \textbf{Vendeur} : Possède les droits pour créer, modifier des annonces et uploader des images (S3).
\end{itemize}

\newpage

%--- 3. SYSTEM SEQUENCE DIAGRAMS ---
\section{Diagrammes de Séquence}

\subsection{Scénario : Messagerie Temps Réel (WebSocket)}
Ce diagramme détaille le flux de communication instantanée implémenté.

\begin{center}
\begin{tikzpicture}
    % Lifelines
    \node[draw, fill=gray!10] (Buyer) at (0,0) {Acheteur (Nav)};
    \node[draw, fill=gray!10] (Server) at (4,0) {Serveur (Spring)};
    \node[draw, fill=gray!10] (DB) at (8,0) {Base de Données};
    \node[draw, fill=gray!10] (Seller) at (12,0) {Vendeur (Nav)};

    \draw[thick] (Buyer) -- (0,-8);
    \draw[thick] (Server) -- (4,-8);
    \draw[thick] (DB) -- (8,-8);
    \draw[thick] (Seller) -- (12,-8);

    % Sequence Interaction
    
    % Step 1
    \node[anchor=west, font=\scriptsize] at (0.1,-1) {1. Connexion STOMP (/ws)};
    \draw[->] (0,-1) -- (4,-1);
    \draw[dashed, ->] (12,-1.2) -- (4,-1.2) node[midway, above, font=\tiny] {Connexion seller};

    % Step 2
    \node[anchor=west, font=\scriptsize] at (0.1,-2) {2. Envoi Message (/app/chat.send)};
    \draw[->] (0,-2) -- (4,-2);

    % Step 3
    
    \draw (4,-2.5) -- (8,-2.5) node[midway, above, font=\tiny] {3. saveMessage()};
    \draw (8,-3) -- (4,-3) node[midway, above, font=\tiny] {MessageSaved};

    % Step 4
    \node[anchor=west, font=\scriptsize, align=left] at (4.1,-4) {4. Broadcast Topic\\(/topic/chatroom/1)};
    \draw[->] (4,-4.5) -- (0,-4.5);
    \draw[->] (4,-4.5) -- (12,-4.5);

    % Step 5
    \node[font=\scriptsize, fill=white, draw] at (6,-6) {Mise à jour UI instantanée via JS};

\end{tikzpicture}
\end{center}

\subsection{Scénario : Création d'Annonce}
Synchronisation entre la base relationnelle et l'index de recherche.

\begin{center}
\begin{tikzpicture}
    % Lifelines
    \node[draw, fill=gray!10] (User) at (0,0) {Utilisateur};
    \node[draw, fill=gray!10] (Ctrl) at (3,0) {ProductController};
    \node[draw, fill=gray!10] (Svc) at (6,0) {ProductService};
    \node[draw, fill=gray!10] (Repo) at (9,0) {MySQL};
    \node[draw, fill=gray!10] (ES) at (12,0) {Elasticsearch};

    \draw[thick] (User) -- (0,-7);
    \draw[thick] (Ctrl) -- (3,-7);
    \draw[thick] (Svc) -- (6,-7);
    \draw[thick] (Repo) -- (9,-7);
    \draw[thick] (ES) -- (12,-7);

    % Scenario
    \draw[->] (0,-1) -- (3,-1) node[midway, above, font=\tiny] {POST /products/new};
    \draw[->] (3,-1.5) -- (6,-1.5) node[midway, above, font=\tiny] {create()};
    
    \draw[->] (6,-2.5) -- (9,-2.5) node[midway, above, font=\tiny] {save(Product)};
    \draw[dashed, ->] (9,-3) -- (6,-3);
    
    \draw[->] (6,-4) -- (12,-4) node[midway, above, font=\tiny] {index(ProductDocument)};
    \draw[dashed, ->] (12,-4.5) -- (6,-4.5);
    
    \draw[dashed, ->] (6,-5.5) -- (3,-5.5);
    \draw[dashed, ->] (3,-6) -- (0,-6) node[midway, above, font=\tiny] {Redirection};

\end{tikzpicture}
\end{center}

\newpage

%--- 4. CLASS DIAGRAM ---
\section{Diagramme de Classe}

Le modèle de domaine reflète la structure relationnelle implémentée avec JPA.

\begin{center}
\begin{tikzpicture}[scale=0.7, transform shape]
    % Classes
    \node[class] (User) {
        \textbf{User}
        \nodepart{second}
        + id: Long\\
        + username: String\\
        + email: String\\
        + password: String\\
        + role: Role
    };

    \node[class, right=of User, xshift=2cm] (Product) {
        \textbf{Product}
        \nodepart{second}
        + id: Long\\
        + title: String\\
        + price: BigDecimal\\
        + status: Status\\
        + location: String
    };

    \node[class, below=of Product, yshift=-1cm] (Image) {
        \textbf{ProductImage}
        \nodepart{second}
        + id: Long\\
        + imageUrl: String\\
        + isPrimary: Boolean
    };

    \node[class, above=of Product, yshift=1cm] (Category) {
        \textbf{Category}
        \nodepart{second}
        + id: Long\\
        + name: String\\
        + icon: String
    };

    \node[class, below=of User, yshift=-1cm] (ChatRoom) {
        \textbf{ChatRoom}
        \nodepart{second}
        + id: Long\\
        + lastMessage: String\\
        + unreadCount: Integer
    };

    \node[class, below=of ChatRoom, yshift=-1cm] (Message) {
        \textbf{ChatMessage}
        \nodepart{second}
        + id: Long\\
        + content: Text\\
        + timestamp: DateTime
    };

    % Relations
    \draw[-{Diamond}] (User) -- node[above] {1} node[below] {*} (Product);
    \draw[-{Diamond}] (Category) -- node[right] {1} (Product);
    \draw[-{Diamond[fill=black]}] (Product) -- node[right] {1} node[left] {*} (Image);
    
    \draw[->] (ChatRoom) -- node[right] {concerne} (Product);
    \draw[->] (ChatRoom) -- node[left] {acheteur} (User);
    \draw[-{Diamond[fill=black]}] (ChatRoom) -- (Message);
    \draw[->] (Message) -- node[left] {expéditeur} (User);

\end{tikzpicture}
\end{center}

\newpage

%--- 5. ARCHITECTURE ---
\section{Architecture Logicielle}

\subsection{Architecture en Couches}
L'application respecte strictement le pattern \textbf{Layered Architecture} :
\begin{itemize}
    \item \textbf{Presentation Layer} : Thymeleaf (SSR) pour le web, JSON pour l'API Mobile.
    \item \textbf{Business Layer} : Services Spring transactionnels encapsulant la logique métier (`@Service`).
    \item \textbf{Persistence Layer} : Interfaces `JpaRepository` pour MySQL et `ElasticsearchRepository`.
    \item \textbf{Database Layer} : MySQL 8.0 (Données relationnelles) et Elasticsearch (Indexation).
\end{itemize}

\subsection{Diagramme de Paquetages (Spring Modulith)}
L'architecture modulaire permet de découpler les fonctionnalités fonctionnelles (User, Product, Chat, Search).

\begin{center}
\begin{tikzpicture}
    % Packages
    \node[package] (UserPkg) at (0, 0) {\textbf{com.marketplace.user}};
    \node[package] (ProdPkg) at (5, 0) {\textbf{com.marketplace.product}};
    \node[package] (ChatPkg) at (0, -4) {\textbf{com.marketplace.chat}};
    \node[package] (SearchPkg) at (5, -4) {\textbf{com.marketplace.search}};
    
    % Dependencies
    \draw[->, thick, dashed] (ProdPkg) -- node[midway, above] {depend} (UserPkg);
    \draw[->, thick, dashed] (ChatPkg) -- node[midway, left] {depend} (UserPkg);
    \draw[->, thick, dashed] (ChatPkg) -- (ProdPkg);
    \draw[->, thick, dashed] (ProdPkg) -- (SearchPkg);

    % Description
    \node[align=left, font=\small] at (10, -2) {
        \textbf{User} : Auth, Profils.\\
        \textbf{Product} : Annonces, Images.\\
        \textbf{Chat} : WebSocket, Salles.\\
        \textbf{Search} : Indexation ES.
    };
\end{tikzpicture}
\end{center}

\section{Architecture de Déploiement}

Le déploiement utilise Docker pour conteneuriser chaque service.

\begin{center}
\begin{tikzpicture}
    % Nodes
    \node[draw, cube, minimum size=2cm, fill=white] (App) at (0,0) {Spring Boot\\(Port 8080)};
    \node[draw, cube, minimum size=1.5cm, fill=blue!10] (MySQL) at (5,2) {MySQL\\(Port 3306)};
    \node[draw, cube, minimum size=1.5cm, fill=yellow!10] (ES) at (5,0) {Elasticsearch\\(Port 9200)};
    \node[draw, cube, minimum size=1.5cm, fill=green!10] (S3) at (5,-2) {LocalStack S3\\(Port 4566)};
    
    \node[draw, cloud, minimum width=3cm, minimum height=1.5cm, fill=gray!5] (Client) at (-5,0) {Navigateur /\\Mobile};

    % Connections
    \draw[->, thick] (Client) -- node[above] {HTTP/WS} (App);
    \draw[->, thick] (App) -- node[above] {JDBC} (MySQL);
    \draw[->, thick] (App) -- node[above] {REST} (ES);
    \draw[->, thick] (App) -- node[below] {AWS SDK} (S3);

    % Boundary
    \draw[dashed, red] (-2, -3.5) rectangle (7, 3.5);
    \node[red] at (2.5, 3.7) {Docker Host};

\end{tikzpicture}
\end{center}

\newpage

%--- 7. IMPLEMENTATION & 8. TESTING ---
\section{Détails d'Implémentation}

\subsection{Stack Technologique}
\begin{itemize}
    \item \textbf{Backend} : Java 17, Spring Boot 3.3.0.
    \item \textbf{Frontend} : Thymeleaf, TailwindCSS, Alpine.js (pour la réactivité du Chat).
    \item \textbf{Données} : MySQL, Elasticsearch, AWS S3 (Images).
    \item \textbf{Protocole} : STOMP over WebSocket pour la messagerie instantanée.
\end{itemize}

\subsection{Fonctionnalités Clés Implémentées}
\begin{enumerate}
    \item \textbf{Messagerie Temps Réel} : Système complet de chat (Salles, Messages, Notifications non-lues) sans rechargement de page.
    \item \textbf{Recherche Hybride} : Utilisation d'Elasticsearch pour des recherches full-text performantes sur les produits.
    \item \textbf{Gestion des Images} : Upload multiple et stockage optimisé sur S3.
    \item \textbf{Authentification} : Sécurisation complète avec Spring Security et BCrypt.
\end{enumerate}

\section{Tests et Qualité}
Une stratégie de test pyramide a été adoptée :
\begin{itemize}
    \item \textbf{Tests Unitaires} : Utilisation de JUnit et Mockito pour isoler la logique métier (`ChatRoomService`, `ProductService`).
    \item \textbf{Tests d'Intégration} : Tests `@SpringBootTest` avec base de données H2 en mémoire pour valider les flux JPA et REST.
    \item \textbf{Vérification Manuelle (E2E)} : Scénarios complets testés via navigateur (Flux Achat/Vente, Chat temps réel) validés par captures d'écran et enregistrements.
\end{itemize}

\section{Conclusion}
Le projet \textit{Oudekenis} a permis de réaliser une plateforme robuste et évolutive. L'adoption de l'architecture \textbf{Spring Modulith} et des \textbf{WebSockets} démontre une maîtrise des technologies modernes du web.
\\
\\
\textbf{Perspectives d'amélioration :}
\begin{itemize}
    \item Ajout d'une application mobile Flutter native (API REST déjà prête).
    \item Mise en place d'un cache Redis pour les sessions distribuées.
    \item Intégration de paiement en ligne.
\end{itemize}

\end{document}
